
\documentclass[11pt]{article}

\usepackage{amsmath}
\usepackage{amssymb}
\usepackage{graphicx}
\usepackage{hyperref}
\usepackage{geometry}
\geometry{margin=1in}

\title{The Prime Exponent Tree (PET): \\
A Structural View of Integer Factorization}

\author{Giancarlo Comneno \\
\small Independent Research}

\date{\today}

\begin{document}

\maketitle

\begin{abstract}
We introduce the \textbf{Prime Exponent Tree (PET)}, a recursive tree representation
of integers derived from prime factorization. In this model, every integer is
represented as a tree whose nodes correspond to prime bases and whose children
recursively encode the structure of the corresponding exponents.

Using a complete dataset of all integers up to $10^6$, we empirically explore the
space of PET structures. Surprisingly, despite the large number of integers,
only a small number of distinct structural shapes appear. Among the first one
million integers, we observe only 78 distinct structural shapes, and the
distribution of these shapes is highly concentrated: seven shapes account for
approximately 87\% of all integers.

We further measure the structural entropy of the PET distribution and observe
that the complexity of integer multiplicative structure grows approximately as
$\log\log N$. These results suggest that the multiplicative structure of integers
is dramatically simpler than their numerical magnitude might suggest.
\end{abstract}

\section{Introduction}

Integers grow rapidly, but their multiplicative structure evolves much more
slowly. Classical results in analytic number theory such as the
Hardy--Ramanujan theorem and the Erd\H{o}s--Kac theorem show that the number
of distinct prime factors of an integer grows on average like $\log\log n$.

However, traditional approaches to integer factorization focus primarily on the
values of primes rather than on the \emph{structure} of factorizations.
In this work we explore a complementary viewpoint: representing integers as
recursive combinatorial objects.

We introduce the \textbf{Prime Exponent Tree (PET)}, a canonical representation
of integers in which prime bases appear as nodes and the structure of exponents
is recursively encoded as subtrees.

Our central empirical question is:

\begin{quote}
How many distinct multiplicative \emph{structures} appear among the integers
up to $N$?
\end{quote}

Using a complete dataset of all integers up to one million, we show that the
space of such structures is surprisingly small and highly organized.

\section{The Prime Exponent Tree}

\subsection{Definition}

Let $N \ge 2$ be an integer with prime factorization

\begin{equation}
N = \prod_i p_i^{e_i}.
\end{equation}

The \textbf{Prime Exponent Tree (PET)} of $N$ is defined recursively as follows:

\begin{itemize}
\item Each prime $p_i$ corresponds to a node.
\item If the exponent $e_i = 1$, the node is a leaf.
\item If $e_i > 1$, the exponent itself is represented by a PET subtree.
\end{itemize}

Thus exponents are treated as integers whose own multiplicative structure
is recursively expanded.

\subsection{Example}

Consider

\begin{equation}
12 = 2^2 \cdot 3.
\end{equation}

The exponent $2$ of the prime $2$ is itself an integer and therefore has
its own PET representation. The resulting structure can be visualized as

\begin{verbatim}
•
├─2
│  •
│  └─2
└─3
\end{verbatim}

This representation captures not only the prime bases but also the recursive
structure of the exponents.

\subsection{Structural Properties}

The PET representation has several useful properties:

\begin{itemize}
\item \textbf{Canonical}: each integer corresponds to a unique PET structure.
\item \textbf{Invertible}: the original integer can be reconstructed exactly.
\item \textbf{Lossless}: no information from the prime factorization is lost.
\end{itemize}

\subsection{Structural Metrics}

Several structural metrics can be defined on PET trees:

\begin{itemize}
\item \textbf{Node count}: total number of nodes in the tree.
\item \textbf{Leaf count}: number of prime factors.
\item \textbf{Height}: maximum recursion depth.
\item \textbf{Branching factor}: number of children of the root.
\end{itemize}

These metrics allow us to study the combinatorial structure of integer
factorizations.

\section{Dataset and Experimental Setup}

We generated a dataset containing the PET representation of all integers
from $2$ to $1,000,000$.

For each integer the dataset stores:

\begin{itemize}
\item the integer value $n$
\item its PET representation
\item structural metrics derived from the tree
\end{itemize}

The resulting dataset contains 999,999 records and occupies approximately
180 MB in JSONL format.

\section{Structural Shape Distribution}

(Section to be completed in the next draft.)

\section{Entropy of Integer Structure}

(Section to be completed.)

\section{Discussion}

(To be completed.)

\section{Conclusion}

(To be completed.)

\end{document}
